% Acronymes du mémoire. À compléter et vérifier.
\newacronym{api}{API}{Application Programming Interface}
\newacronym{clahe}{CLAHE}{Contrast Limited Adaptive Histogram Equalization}
\newacronym{fps}{FPS}{Frames Per Second}
\newacronym{hsv}{HSV}{Hue, Saturation, Value}
\newacronym{ia}{IA}{Intelligence Artificielle}
\newacronym{iot}{IoT}{Internet of Things}
\newacronym{json}{JSON}{JavaScript Object Notation}
\newacronym{lan}{LAN}{Local Area Network}
\newacronym{mfa}{MFA}{authentification multifacteur}
\newacronym{mjpeg}{MJPEG}{Motion JPEG}
\newacronym{rest}{REST}{Representational State Transfer}
\newacronym{rgpd}{RGPD}{Règlement général sur la protection des données}
\newacronym{sql}{SQL}{Structured Query Language}
\newacronym{ups}{UPS}{Uninterruptible Power Supply}
\newacronym{yolo}{YOLO}{You Only Look Once}
\newacronym{far}{FAR}{taux de fausse acceptation}
\newacronym{frr}{FRR}{taux de faux rejet}
\newacronym{lfw}{LFW}{Labeled Faces in the Wild}
\newacronym{apdp}{APDP}{Autorité de protection des données à caractère personnel}
\newacronym{ftc}{FTC}{Federal Trade Commission}
\newacronym{uml}{UML}{Unified Modeling Language}
\newacronym{csv}{CSV}{Comma-Separated Values}
\newacronym{rtsp}{RTSP}{Real-Time Streaming Protocol}
\newacronym{https}{HTTPS}{HyperText Transfer Protocol Secure}
\newacronym{lcd}{LCD}{Liquid Crystal Display}
\newacronym{pin}{PIN}{numéro d'identification personnel}
\newacronym{sms}{SMS}{Short Message Service}
\newacronym{gsm}{GSM}{Global System for Mobile Communications}
\newacronym{roc}{ROC}{Receiver Operating Characteristic}
\newacronym{cnil}{CNIL}{Commission nationale de l'informatique et des libertés}
\newacronym{nist}{NIST}{National Institute of Standards and Technology}
\newacronym{coco}{COCO}{Common Objects in Context}
\newacronym{usb}{USB}{Universal Serial Bus}

% Glossaire : définitions générales, à valider et ajuster par l'étudiant.
\newglossaryentry{antispoofing}{name={Anti-usurpation (anti-spoofing)}, description={Ensemble des techniques qui visent à empêcher qu'un système biométrique soit trompé par une reproduction du visage (photographie imprimée, écran, masque)}}
\newglossaryentry{cloudcomputing}{name={Cloud Computing}, description={Informatique en nuage : modèle dans lequel les données sont traitées et stockées sur des serveurs distants, accessibles par Internet}}
\newglossaryentry{edge}{name={Edge Computing}, description={Modèle d'architecture dans lequel le traitement des données est effectué localement, à proximité de la source, plutôt que sur un serveur distant}}
\newglossaryentry{embedding}{name={Embedding}, description={Représentation numérique, sous forme de vecteur, d'un objet (ici un visage) permettant de mesurer sa ressemblance avec un autre}}
\newglossaryentry{liveness}{name={Liveness (test de vivacité)}, description={Vérification que la personne présentée devant la caméra est bien vivante, par exemple en lui demandant de sourire ou de cligner des yeux}}
\newglossaryentry{trustscore}{name={Score de confiance (Trust Score)}, description={Indicateur qui augmente à mesure que les étapes d'authentification sont validées successivement}}

\newglossaryentry{secureenclave}{name={Secure Enclave}, description={Composant sécurisé des appareils Apple qui protège les données du visage et réalise la comparaison}}

\newglossaryentry{enrolement}{name={Enrôlement}, description={Enregistrement d'une personne dans le système : son visage est photographié et transformé en vecteurs, rangés dans un profil}}
\newglossaryentry{gabarit}{name={Gabarit}, description={Ensemble des données biométriques conservées pour reconnaître une personne (dans BioGate, des vecteurs de nombres)}}
\newglossaryentry{distancecosinus}{name={Distance cosinus}, description={Mesure de la différence entre deux vecteurs : elle vaut 0 pour deux vecteurs identiques et augmente quand ils diffèrent}}
\newglossaryentry{fenetregrace}{name={Fenêtre de grâce}, description={Délai laissé, après la détection d'un inconnu, pour reconnaître une personne autorisée avant de confirmer l'intrus et d'envoyer l'alerte}}
\newglossaryentry{jeton}{name={Jeton}, description={Code secret remis par le serveur au tableau de bord après le couplage, et joint à chaque demande pour prouver que le tableau de bord est autorisé}}
\newglossaryentry{route}{name={Route}, description={Adresse du serveur à laquelle un programme adresse une demande, par exemple la liste des profils}}
\newglossaryentry{opencv}{name={OpenCV}, description={bibliothèque logicielle libre de traitement d'images et de vision par ordinateur}}
\newglossaryentry{fastapi}{name={FastAPI}, description={cadre de développement en Python pour construire des interfaces de programmation accessibles par le réseau}}
\newglossaryentry{python}{name={Python}, description={langage de programmation}}
\newglossaryentry{javascript}{name={JavaScript}, description={langage de programmation des pages web}}
\newglossaryentry{electron}{name={Electron}, description={cadre de développement qui permet de créer des applications de bureau avec les technologies du web}}
\newglossaryentry{react}{name={React}, description={bibliothèque JavaScript pour construire des interfaces utilisateur}}
\newglossaryentry{sqlite}{name={SQLite}, description={système de base de données contenu dans un seul fichier, sans serveur}}
\newglossaryentry{tensorflow}{name={TensorFlow}, description={bibliothèque d'apprentissage automatique développée par Google}}
\newglossaryentry{mediapipe}{name={MediaPipe}, description={bibliothèque de Google qui propose des modèles de détection de visages et de points du visage}}
\newglossaryentry{vggface}{name={VGG-Face}, description={modèle de reconnaissance faciale par réseau de neurones, développé à l'université d'Oxford}}
\newglossaryentry{arduino}{name={Arduino}, description={environnement de programmation de cartes électroniques}}
\newglossaryentry{bluetooth}{name={Bluetooth}, description={liaison radio à courte portée entre appareils}}
\newglossaryentry{wifi}{name={Wi-Fi}, description={réseau local sans fil}}
\newglossaryentry{telegram}{name={Telegram}, description={application de messagerie instantanée}}
\newglossaryentry{wikicommons}{name={Wikimedia Commons}, description={base en ligne de photographies et de fichiers libres de droits, gérée par la fondation Wikimedia}}
\newglossaryentry{raspberrypi}{name={Raspberry Pi}, description={micro-ordinateur monocarte peu coûteux}}
\newacronym{html}{HTML}{HyperText Markup Language}
\newacronym{css}{CSS}{Cascading Style Sheets}
\newglossaryentry{faceid}{name={Face ID}, description={système de reconnaissance faciale d'Apple pour déverrouiller ses appareils}}
\newglossaryentry{googlenest}{name={Google Nest}, description={gamme de caméras et d'appareils connectés de Google}}
\newglossaryentry{frigate}{name={Frigate}, description={logiciel libre de vidéosurveillance qui s'exécute en local et détecte des objets dans les images}}
